\chapter{Conclusion and Future Work}

\section{Conclusion}

This project presented the design and implementation of a User-Friendly Autonomous Parcel Delivery System with Adaptive Navigation aimed at addressing the growing demand for efficient and intelligent indoor logistics solutions. The system was developed entirely in a simulation environment using ROS~2 and Gazebo, focusing on reliable navigation, intelligent path planning, and smooth trajectory execution so that robot does not have to slow down during turns and effiently aviod any walking person.

The primary objective of the project was to design a mobile robot capable of navigating structured indoor environments such as university buildings and office corridors while safely avoiding obstacles and delivering parcels to specified goal locations.It can also be placed in special enivornemnts like isolation wards or hospital to avoid virus transmission. This objective was successfully achieved through the integration of an Improved A* (IA*) algorithm for global path planning, cubic spline interpolation for trajectory smoothing, and a Model Predictive Control (MPC) strategy for accurate motion tracking.

Different path planning algorithms were studies but he Improved A* algorithm proved effective in generating collision-free and near-optimal paths on grid maps. The addition of obstacle inflation and waypoint pruning enhanced both safety and computational efficiency. Since grid-based planners naturally produce piecewise-linear paths with sharp turns, spline-based smoothing was applied to generate continuous and kinematically feasible trajectories suitable for differential-drive robotsso so that does not have to slow down during curves and effiently aviod any walking person. The MPC controller ensured stable and accurate tracking of these trajectories across different simulated environments.

All the system was deployed in the context of the ROS 2 with the use of modular nodes and traditional communication patterns. A web-based control interface was also designed simultaneously to allow users to define delivery goals, simulate the spatial environment, and track the actions of the robot in real time. Extensive testing also revealed that the robot is able to navigate complicated indoor setups, avoid obstacles and get to the predefined target points with a good degree of consistency.

On the whole, it can be concluded that the project has achieved its desired goals and supported the fact that simulation-based development is a powerful and safe approach to the development of autonomous delivery systems. The article provides a strong background of further research and practical implementation of smart indoor logistics robots.

\section{Future Work}
 As much as the proposed system has demonstrated good performance in the simulated indoor settings, there are yet to be explored promising research directions. Those extensions might increase the intelligence, scalability and usefulness of the system in real world significantly.

\begin{itemize}
    \item \textbf{Real-World Deployment and Sensor Fusion:}  
 One of the major follow-up tasks is to transfer the system to a real robot. It implies connecting real sensors LiDAR, RGB -D cameras, IMUs, wheel encoders, etc. With this advanced sensor fusion, e.g. Extended Kalman Filters or factor-graph SLAM, the robot should be able to more accurately localise itself in the presence of all the real-world noises and uncertainty.
 
    \item \textbf{Learning-Based Navigation and Decision Making:}  
  As much as the proposed system has demonstrated good performance in the simulated indoor settings, there are yet to be explored promising research directions. Those extensions might increase the intelligence, scalability and usefulness of the system in real world significantly.
  
  The next generation implementation may incorporate reinforcement learning or imitation learning to enable the robot to gain navigation policies by experience. That would allow it to be more adaptive to complex and crowded environments.
    \item \textbf{Multi-Robot Coordination:}  
   We can totally scale it, Hey, we could run a number of delivery robots, at a time. It would then be necessary to determine how to divide work among them, establish a manner of them communicating with each other, and coordinate their pathways. It would make the entire process much more efficient in large institutions such as hospitals or college campuses.
   
  
    \item \textbf{Semantic Mapping and Context-Aware Navigation:}  
  
  The inclusion of semantic SLAM would enable the robot to actually see things such as doors, elevators, people, furniture so it can tag them. Knowing that, the robot may make more intelligent and higher level decisions and maintain itself around people safer.
  
    \item \textbf{Dynamic Global Replanning and Prediction:}  
 
 We can also predictively global replan in future, so the robot will actually guess where people are heading and dynamically adjust its route on the fly according to the changing environment.
 

    \item \textbf{Human--Robot Interaction (HRI):}  
   Lastly, in case we add voice commands, phone-to-phone push notifications, and certain visual feedback, the system would be much more friendly and more acceptable in daily environments.
   
    \item \textbf{Energy-Aware and Sustainable Navigation:}  
     Autonomous charging plans can be combined with energy-efficient path planning to allow the robotic systems to operate sustainably and over a long period.
    
    \item \textbf{Cloud and IoT Integration:}  
   The integration of the system with cloud services as well as the IoT infrastructure would allow tracking the fleet, analyzing the data in real-time, and improving it on the basis of the actual data related to the real operation.
\end{itemize}
